\documentclass[11pt,twoside]{book}
\usepackage[utf8]{inputenc}
\usepackage[spanish,es-nodecimaldot]{babel}
\usepackage[papersize={5.5in,8.5in}, margin=0.75in, bindingoffset=0.25in]{geometry}
\usepackage{setspace}
\onehalfspacing % Un poco de aire para legibilidad
\usepackage{charter} % Una fuente más clásica y legible para libros
\usepackage{bookmark} % Mejora la creación del índice en el PDF
\usepackage{hyperref}

\title{\textbf{El Velo Gravitacional}}
\author{Carlos G. Ramirez}
\date{} % Sin fecha en la portada para mayor limpieza

\begin{document}

\maketitle

\frontmatter
\tableofcontents

\mainmatter

\chapter*{Prólogo: EL REGALO}
\addcontentsline{toc}{chapter}{Prólogo: EL REGALO}
\markboth{Prólogo}{Prólogo}

\section*{3 de febrero de 2114}

Ethan Ramírez despertó con el sabor amargo y seco de la calima industrial pegado a la lengua. El aire del Distrito RD no era solo aire; era una suspensión pesada cargada de residuos de procesado y humedad ácida, donde la dispersión de Rayleigh se rendía ante una bruma permanente que convertía el sol en un disco amarillento, una pupila enferma que observaba desde el otro lado de la polución. Diecisiete años. Hoy cumplía diecisiete años, una cifra que en la periferia no se medía en pasteles ni deseos, sino en la capacidad de los pulmones para seguir filtrando el óxido de la existencia cotidiana.

Se incorporó en su pequeño catre, que chirrió en protesta. Se frotó los ojos, tratando de quitarse el cansancio que parecía haberse filtrado en sus huesos durante la noche. A través de la ventana rota de su habitación, sellada parcialmente con cinta plástica que vibraba con el viento, podía ver el brillo insultante del domo de Goglee City a lo lejos. Era una semiesfera perfecta de cristal y nanotecnología, una luna artificial cautiva que bañaba a los privilegiados en una luz blanca y pura, recordándole a cada segundo que había un mundo de aire filtrado y cielos azules al que él nunca tendría invitación.

---¡Ethan! ---La voz de Carlos, profunda y rasposa como el papel de lija, retumbó desde el taller de abajo---. ¡Despierta, mijo! ¡Baja, que lo que te tengo no espera!

Ethan se vistió en silencio, poniéndose la misma camiseta de tirantes gastada y los pantalones de carga manchados de grasa que eran casi una segunda piel. Al final de la escalera de metal, el calor del taller ya empezaba a subir, trayendo consigo el aroma reconfortante del aceite de motor y el ozono de las soldaduras.

El taller de Carlos era su santuario. Era un espacio angosto, encajado entre edificios condenados, donde el techo goteaba cuando la lluvia ácida caía con fuerza, pero para Ethan era el centro del universo. Entre pilas de transmisiones rotas y cables de fibra óptica deshilachados, su padre obraba milagros. Carlos le había enseñado que una pieza de metal no estaba muerta mientras pudieras sentir su vibración. ``Escucha el motor, Ethan'', solía decirle, ``él te dirá dónde le duele antes de que se rompa''.

Carlos lo esperaba en medio del caos organizado de herramientas. A sus cincuenta y seis años, sus manos eran un mapa de cicatrices y quemaduras, pero sus ojos oscuros conservaban un brillo de picardía que el Distrito no había logrado apagar. Carlos era un hombre forjado en la lógica del metal; para él, el universo era una máquina compleja regida por leyes físicas que no necesitaban oraciones para cumplirse. Aunque usaba la bendición por puro hábito cultural, su verdadera fe residía en la precisión de un motor bien alineado.

Detrás de él, envuelta en una lona de polímero manchada de mil batallas mecánicas, se erguía una silueta inconfundible. Una promesa de velocidad grabada en sus líneas ocultas.

---Cion papi. Papi!... ---Ethan se detuvo, sintiendo un nudo en la garganta---. ¿No me digas que...?

Carlos sonrió. ---¡Jajaja! Dios te bendiga, mijo. ---Era una expresión de orgullo que le restó diez años de encima.

---Feliz cumple, mijo. No es una de esas máquinas brillantes que usan los \textit{popis} de la ciudad alta, pero tiene un núcleo de inducción gravitacional Clase 3 que rescaté de un carguero estrellado. Es puro hierro, Ethan. Del bueno.

Ethan se acercó como si temiera que la visión se desvaneciera. Al retirar la lona, el polvo bailó en la luz mortecina. El chasis estaba rayado y el asiento de cuero sintético tenía parches, pero bajo la carcasa de aleación, Ethan podía sentir el pulso de algo poderoso. Era el regalo más valioso que alguien en el Distrito RD podía recibir: un par de alas de metal.

---Es... es increíble, papi ---susurró Ethan, pasando los dedos por el manillar. Notó la resistencia familiar de los controles, la firmeza de la horquilla---. ¿Cómo conseguiste las piezas para el inversor G?

---Eso es secreto profesional, mijo ---rio Carlos, aunque luego tosió brevemente, ocultando la boca---. La he estado armando por piezas estas últimas noches. No levita todavía, el sistema de fase está muerto, pero el motor ruge como un toro. El resto... el resto lo vas a poner tú. Es tu proyecto ahora.

Carlos suspiró, clavando la vista en un viejo monitor que mostraba telemetría astronómica llena de estática. Sus ojos, antes brillantes, se sombrearon por un instante.

---Hay algo en esa data que no me cuadra, Ethan. No me gusta. ---Carlos señaló el monitor, donde una línea de telemetría oscilaba con una arritmia inquietante---. La Cúpula dice que es simple ruido, pero un mecánico sabe cuando una máquina va a reventar aunque el panel diga que todo está en verde. Por otro lado hijo, aprende a sentir el espacio-tiempo si tanto te gustan las motowaves. Puede que pronto necesites saber por dónde escapar cuando el cielo decida cerrarse.

Ethan cerró los ojos y, por un momento, no vio las paredes desconchadas del taller. Vio las historias que su padre le contaba por las noches, cuando el ruido de la ciudad se calmaba. Historias de una isla llamada Quisqueya, donde el mar no era negro por el petróleo sino de un azul tan intenso que dolía mirarlo. Donde la gente no necesitaba máscaras para respirar y el sol te besaba la piel en lugar de quemarla.

Esa herencia latía en él. Era lo que lo mantenía en pie. Se giró hacia su padre, buscando ese vínculo que era su única ancla.

---Papi... Ción.

Ethan inclinó levemente la cabeza, un gesto de respeto antiguo que Carlos exigía no por ego, sino por memoria. Era el cordón umbilical que los unía a sus ancestros, a los que cruzaron el mar antes de que el mar se lo tragara todo.

Carlos puso su mano callosa sobre la cabeza de Ethan. El contacto fue cálido, cargado de una bendición que no necesitaba de iglesias para ser sagrada.

---Dios me lo bendiga y me lo haga un hombre de bien, mijo. Y que nunca se le olvide que, aunque estemos en el fango, nosotros somos gente de cielo.

Ethan asintió, sintiendo el peso de la responsabilidad. Miró su moto de nuevo. No era solo chatarra recuperada. Era el inicio de algo.

---La llamaré \textit{El Fénix} ---dijo con firmeza---. Porque va a salir de este agujero volando.

Carlos lo miró en silencio, con una mezcla de esperanza y un miedo que Ethan aún no alcanzaba a comprender. Sabía que al darle esa moto, le estaba dando un sueño. Y en el Distrito RD, los sueños solían tener un precio muy alto.

Pero hoy no era día para pensar en precios. Hoy, Ethan Ramírez tenía diecisiete años y, por primera vez en su vida, sentía que la gravedad ya no pesaba tanto.

\bigskip
\begin{center}***\end{center}
\bigskip

\textbf{FIN DEL PRÓLOGO}


\bigskip
\begin{center}***\end{center}
\bigskip


\part{EL APRENDIZAJE}

\chapter{PRIMEROS PASOS}

\section*{Febrero de 2114 - Una semana después del cumpleaños}

Ethan despertó con el sonido rítmico y seco de metal golpeando metal. El sol del Distrito RD, una mancha pálida y difusa, apenas lograba teñir el smog persistente de un tono naranja sucio que recordaba al óxido. Para cualquier otro, ese amanecer sería deprimente, pero para Ethan, el color del cielo hoy era lo de menos. Hoy, en el centro de su habitación mental, solo existía una cosa: la motowave.

Bajó las escaleras de metal de dos en dos, sintiendo la vibración del taller en las plantas de sus pies. Abajo, Carlos ya estaba sumergido en su elemento. Las chispas de la soldadura saltaban en arcos azulados, iluminando momentáneamente las paredes cubiertas de herramientas y esquemas técnicos amarillentos.

---Buenos días, mijo ---dijo Carlos sin apartar la vista de la pieza que unía---. Te me adelantaste al despertador.

---¿Cómo voy a dormir, papi? ---Ethan se acercó a la motowave, pasando los dedos por el chasis. La sensación del metal frío y áspero le recorrió el brazo---. Anoche soñé con el esquema del motor de inducción. Tenemos trabajo que hacer.

Carlos apagó el soplete y se levantó la careta protectora. Sus ojos reflejaban una mezcla de ternura y la severidad propia del oficio.

---Antes de que te llenes de grasa hasta las orejas, hay que ser realistas. ---Carlos golpeó el metal con los nudillos---. El chasis está sólido, pero el sistema eléctrico es un nido de ratas. Los propulsores gravitacionales nunca se instalaron y el núcleo necesita una alineación de fase que nos va a sacar canas verdes. Si piensas volar con esto, Ethan, vas a necesitar más que buena voluntad.

Ethan sacó su cuaderno de notas, cuyas páginas estaban hinchadas por la humedad y manchadas de grafito. Había diagramas, cálculos de fuerza G y bocetos de aerodinámica que había copiado de manuales desechados.

---Empezaremos por el cableado principal ---dijo Ethan, con una determinación que hizo que Carlos asintiera en silencio---. Necesito ver si el bus de datos todavía responde.

\bigskip
\begin{center}***\end{center}
\bigskip

Tres horas más tarde, Ethan estaba enfrascado en una batalla contra una maraña de cables de cobre oxidados. El taller apestaba a decapante y sudor, pero Ethan apenas lo notaba. Estaba buscando la ``chispa'' de la máquina, ese eco eléctrico que le confirmara que \textit{El Fénix} estaba vivo.

---¡Ethan! ¡Klk manín! ¿Estás metido en el hierro todavía?

Esa voz, cargada de una energía eléctrica que rivalizaba con los condensadores del taller, solo podía ser de Rafa. Su mejor amigo cruzó el umbral moviendo los hombros al ritmo de una música imperceptible, cargando una mochila que traqueteaba con el sonido de componentes electrónicos.

---¡Aquí atrás, Rafa! ---respondió Ethan, emergiendo de entre las piezas con una mancha de aceite en la frente.

---¡Mierda, manito, pero tú te has vuelto un topo! ---exclamó Rafa, dándole un choque de hombros---. Escuché que el viejo Carlos se botó y te soltó la motowave. No está mal para ser chatarra del Distrito, ¿eh?

---No la llames chatarra, loco ---dijo Ethan con una sonrisa, aunque su tono dominicano marcaba el respeto por la máquina---. Es un proyecto. Y uno grande.

---Relájate, hermano, yo sé que para ti esto es oro. ---Rafa dejó caer su mochila y se asomó al interior del motor---. Pero deja que el experto le eche un ojo a ese sistema eléctrico. Eso parece un nido de serpientes muertas de Goglee City. Yo te ayudo con la electrónica, que tú sabes que en eso soy el final.

Ethan se rió, agradecido por la oferta. Rafa era un genio improvisando circuitos con piezas de desecho.

\bigskip
\begin{center}***\end{center}
\bigskip

Al mediodía, el calor se volvió denso y el aire en el taller se estancó. Yeisly llegó como un torbellino, cargando una bolsa que emanaba el aroma celestial del arroz con habichuelas.

---¡Mierda loco, qué jebi está esa vaina! ---soltó Yeisly nada más ver la moto sobre el caballete---. ¡Dime que no es verdad, primo! ¿De verdad vamos a poner a volar este hierro?

---Esa es la idea, Yeisly ---respondió Ethan, recibiendo la comida---. Pero todavía falta mucho.

Yeisly rodeó la moto con los ojos entrecerrados, como si estuviera encuadrando una pintura. Ella veía colores donde los demás veían óxido.

---Tiene líneas potentes. De verdad. Cuando esté lista, yo le voy a meter un diseño que la gente va a creer que es una nave de la Ciudad Aérea. Quiero que el Fénix se vea desde el domo.

Acompañando a Yeisly estaba Unica. Ella, al contrario que la prima de Ethan, se mantuvo en silencio unos instantes, observando la estructura con una mirada analítica y calmada.

---Ethan, las especificaciones del fabricante para este modelo sugieren que el chasis debe ser reforzado ---comentó Unica con su voz educada y precisa---. Si planeas alcanzar velocidades de carrera, los puntos de anclaje de los propulsores sufrirán una fatiga estructural inaceptable.

---¿Titanio? ---preguntó Ethan, reconociendo la lógica en las palabras de su amiga.

---Exactamente. Placas de refuerzo de titanio o, al menos, una aleación de cromo-molibdeno ---continuó Unica, sin rastro de jerga en su discurso---. El costo en créditos es elevado, pero es una medida de seguridad necesaria si valoras tu integridad física.

Rafa levantó la vista de un mazo de cables que estaba pelando.

---¿Y de dónde vamos a sacar los créditos? Esa vaina no cae del cielo.

---Conseguiremos la chatarra ---dijo Unica, saliendo de debajo de la moto con una pequeña linterna técnica---. Conozco a un proveedor en la Plaza que tiene excedentes de las fábricas corporativas. No será nuevo, pero con un poco de pulido, su integridad estructural será perfecta para lo que necesitamos.

Carlos, que observaba la escena desde la puerta, intervino con voz tranquila pero firme.

---Yo puedo dar unas clases extra de mecánica pesada en el muelle de carga. Entre todos sacaremos esto adelante.

Ethan miró a su alrededor. Estaban todos allí: su padre, su mejor amigo, su prima artista y la mente técnica de Unica. Gente que no tenía nada, pero que estaba dispuesta a darle lo poco que tenían para que él pudiera volar.

---Gracias ---dijo Ethan, sintiendo que la garganta se le cerraba un poco. Su acento dominicano se volvió más suave, más sentido---. No sé cómo agradecerles esto.

---No digas nada, manito ---le cortó Rafa, dándole una palmada en la espalda---. Solo gana esa maldita carrera por nosotros.

---¡Eso mismo! ---añadió Yeisly con una carcajada---. Que si perdemos, vamos a tener que comer aire reciclado por tres meses para pagar los platos rotos.

Esa noche, cuando el silencio finalmente reclamó el taller, Ethan se quedó solo frente a la carcasa desarmada de \textit{El Fénix}. Con el cuaderno sobre las rodillas, miró su boceto. Ya no veía solo una moto vieja; veía el esfuerzo de su gente, la esperanza de un distrito entero que se negaba a ser olvidado. Se recostó contra la pared fría y cerró los ojos, imaginando el viento puro del domo golpeándole el rostro.

---Algún día ---susurró para sí mismo, con la voz llena de una promesa solemne---. Algún día vas a volar.

Desde la sombra del pasillo, Carlos lo miraba. Sentía orgullo, sí, pero también un miedo gélido que le recorría la espalda. Sabía que El Velo Gravitacional no premiaba el talento tanto como castigaba la esperanza. Era una carrera diseñada por los de arriba para recordarles a los de abajo cuál era su lugar.

Cerró los ojos y, en la oscuridad, elevó una plegaria silenciosa a los antiguos dioses de la isla.

\textit{Protégelo. Por favor, que esta vez sea diferente.}

\bigskip
\begin{center}***\end{center}
\bigskip

\textbf{FIN DEL CAPÍTULO 1}


\bigskip
\begin{center}***\end{center}
\bigskip

\chapter{EL CÍRCULO}

\section*{Abril de 2114 - Dos meses después}

La Plaza de Chatarra nunca cerraba sus párpados de neón. Incluso a las tres de la madrugada, el lugar palpitaba con una energía febril y clandestina. Era un laberinto de contenedores de carga apilados, donde el aire sabía a ozono quemado, aceite rancio y el aroma picante de las frituras sintéticas que se vendían en los puestos de las esquinas. En la Plaza, el silencio era una anomalía; siempre había un martillero golpeando metal, el zumbido de un generador ilegal o el murmullo de negociaciones que nunca verían la luz del día.

Ethan caminaba con paso firme, tratando de ignorar las pantallas holográficas que proyectaban anuncios de Goglee City, ahora distorsionados y llenos de estática. A su lado, Rafa se movía como si esquivara golpes invisibles, con los ojos bien abiertos. Habían trabajado como mulas durante ocho semanas, doblando turnos en el muelle y vendiendo cada gramo de cobre que lograban rescatar de los vertederos. Hoy, finalmente, el sueño de la levitación estaba al alcance de sus dedos.

---Oye, manito, chequea esto ---susurró Rafa, señalando su tableta agrietada---. El Tocayo dice que el cargamento de inductores ya llegó. Puesto 47-B, allá por el callejón de los Sifones. Si ese viejo no nos clava el cuchillo con el precio, hoy mismo \textit{El Fénix} despega.

Encontraron el puesto de El Tocayo oculto tras una cortina de cables eléctricos que escupían chispas ocasionales. El viejo comerciante, con una mandíbula de cromo que chirriaba al hablar, estaba sentado sobre un trono de memorias RAM obsoletas.

---Los chicos de la motowave ---gruñó El Tocayo, dejando escapar una nube de vapor químico de sus pulmones artificiales---. ¿Me trajeron el ``cariño'' o vinieron a hacerme perder el tiempo?

Ethan sacó un fajo de créditos arrugados, el fruto de su sudor y sus noches en vela.

---Quinientos créditos, El Tocayo. Lo prometido por el par de propulsores gravitacionales Clase 3.

El Tocayo empezó a contar el dinero con una lentitud exasperante, pero antes de que pudiera terminar, una sombra se proyectó sobre el mostrador.

---Esa mercancía ya tenía nombre, viejo ---dijo una voz que cortó el aire como una cuchilla.

Ethan se giró. Frente a él estaba una chica que parecía haber nacido del mismo fuego de la chatarra. Tenía el pelo de un rosa eléctrico que destacaba contra su piel canela, y llevaba una chaqueta de cuero llena de parches con símbolos de resistencia. Sus ojos tenían un brillo desafiante, una mezcla de rabia y una inteligencia afilada que Ethan nunca había visto.

---Oye, chico, busca tus piezas en otro lado ---dijo ella, con un marcado acento cubano que le daba ritmo a su amenaza---. Esos inductores me pertenecen por reserva de palabra.

---En la Plaza no hay palabras, niña ---soltó El Tocayo con una risa metálica---. Solo hay créditos sobre la mesa. Y estos muchachos pagaron primero.

La chica se acercó a Ethan, invadiendo su espacio personal. Era un poco más alta que él, y su presencia desprendía un olor a pólvora y lluvia ácida. Lo estudió de arriba abajo, deteniéndose en sus manos manchadas de aceite.

---¿Y tú para qué quieres unos Clase 3? ---preguntó ella, suavizando apenas el tono---. Eso es mucha potencia para un simple carro de carga.

---No es para carga ---respondió Ethan, sosteniendo la mirada---. Es para una motowave. Voy a correr en El Velo Gravitacional.

La chica soltó una carcajada seca, carente de alegría.

---¿El Velo? ¿De verdad vas a prestarte para el circo de las corporaciones? ¿Vas a dejar que los popis apuesten por tu vida mientras te matas para su entretenimiento?

---Voy a ganar ---dijo Ethan, y su acento dominicano se hizo más marcado por la tensión---. Y con ese premio, voy a sacar a mi gente de este agujero.

---Nadie gana contra la casa, bobo ---replicó ella, cruzándose de brazos---. El Velo es el opio del Distrito RD. Les dan una migaja de esperanza para que no se levanten y quemen ese domo de cristal.

---Entonces déjame ser la excepción ---propuso Ethan---. Si fallo, seré una estadística más. Si gano, habré demostrado que ellos no son invencibles.

La chica lo miró intensamente, buscando alguna grieta en su resolución. Al no encontrarla, hizo un gesto de desdén con la mano y se apartó.

---Llevátelos, pues. Pero cuando estés ahí arriba y veas la sangre en el asfalto, acuérdate de Sony. Acuérdate de que te lo advertí.

---Sony... ---repitió Ethan---. Gracias.

---No me agradezcas nada, que me debes una ---dijo ella antes de desaparecer entre los puestos, moviéndose con una agilidad felina.

Rafa, que había estado conteniendo el aliento, soltó un silbido largo.

---¡Mierda, manito! Esa jeva es fuego puro. Creo que te acaba de echar un mal de ojo o se enamoró de ti, una de dos.

\bigskip
\begin{center}***\end{center}
\bigskip

Dos semanas más tarde, el taller de Carlos era un nido de tensión. Los nuevos propulsores estaban anclados al chasis de \textit{El Fénix}, brillando con un azul cobalto que iluminaba las caras expectantes del equipo. Instalar inductores gravitacionales en un chasis viejo era como intentar meter el alma de un demonio en un cuerpo de madera; la alineación tenía que ser perfecta o la moto se desgarraría a sí misma en el aire.

Carlos supervisaba cada soldadura, cada ajuste de software que Unica realizaba en la consola integrada.

---La inducción está estable en un 92% ---informó Unica, con su tono pausado y preciso---. Ethan, si elevamos el flujo ahora, el sistema de fase debería estabilizarse.

Ethan subió a la moto. El asiento, ahora reforzado, se sentía firme. Respiró hondo, cerró los ojos un segundo y activó el núcleo.

Un zumbido sordo empezó a vibrar en el taller. Las herramientas sobre las mesas empezaron a tintinear. Lentamente, \textit{El Fénix} se despegó del suelo. Diez centímetros. Veinte. Medio metro. La moto flotaba, desafiando siglos de física, sostenida por la pura voluntad y la tecnología rescatada.

---¡Wepa! ¡Lo logramos! ---gritó Rafa, lanzando su gorra al aire.

---¡Mierda, primo, que eso está volando de verdad! ---exclamó Yeisly, grabando la escena con su terminal.

Ethan sentía la moto vibrar bajo él, como un animal salvaje que empezaba a despertar. La adrenalina le recorría el cuerpo. Ya no era solo un sueño; era una realidad que flotaba frente a sus ojos.

\bigskip
\begin{center}***\end{center}
\bigskip

Esa misma noche, mientras Ethan revisaba los sensores térmicos en el silencio del taller, Sony apareció de nuevo. Entró sin llamar, como si el taller fuera una extensión de su propia casa roja en los barrios bajos.

---Oí que el pájaro de hierro ya tiene alas ---dijo, apoyándose en la pared---. Nada mal para un novato.

---¿Qué haces aquí, Sony? ---preguntó Ethan, aunque en realidad se alegraba de verla.

---Vengo a asegurarme de que no te mates antes de la carrera. Sería un desperdicio de inductores ---respondió ella, acercándose a la moto para inspeccionar el trabajo de Unica---. ¿Sabes lo que es la técnica del ``Surfeo de Geodésicas''?

Ethan negó con la cabeza.

---Es lo que te va a salvar la vida cuando los otros corredores intenten sacarte de la pista con un Martillo Gravitacional ---explicó Sony. Sus ojos brillaron con fervor revolucionario---. Los de Goglee City creen que la gravedad es una ley que ellos poseen. Nosotros, los que vivimos en las sombras, sabemos que es solo una ola que hay que aprender a surfear.

Se sentaron en el suelo del taller, compartiendo una botella de refresco templado. Sony le habló de su infancia en los sectores cubanos de la costa, antes de que el agua subiera tanto que solo quedaron las puntas de los edificios. Le habló de su familia, perdida en la Gran Inundación, y de por qué odiaba tanto el domo.

---Yo no quiero un trofeo, Ethan. Yo quiero que el domo caiga. Quiero que el aire puro sea para todos ---dijo ella, y su voz tembló ligeramente.

Ethan puso su mano sobre la de ella. Sony no la retiró.

---Si gano, Sony, tendré voz. Podré decirles lo que pasa aquí abajo.

---Eres tan ingenuo... y tan terco ---susurró ella, acercándose a él.

El beso fue inesperado y eléctrico. Sabía a rebelión y a la desesperación de dos almas que se habían encontrado en el fin del mundo. Cuando se separaron, Sony le dedicó una sonrisa triste.

---Si vas a hacer esto, Ethan, hazlo bien. No seas un peón más. Sé el fénix que incendia su tablero.

Ethan la observó irse, sintiendo que ahora tenía otra razón más para ganar. No solo por su padre, no solo por Distrito RD. Ahora corría por la chica del pelo rosa que creía que el mundo podía ser quemado y reconstruido.

\bigskip
\begin{center}***\end{center}
\bigskip

\textbf{FIN DEL CAPÍTULO 2}


\bigskip
\begin{center}***\end{center}
\bigskip

\chapter{EL VIEJO TORO}

\section*{Junio de 2114}

\textit{El Fénix} estaba dejando de ser un montón de chatarra apilada para convertirse en una obra de arte funcional. Bajo la luz amarillenta del taller, las líneas de su chasis ahora fluían con una elegancia depredadora. Yeisly había pasado noches enteras trabajando en el diseño de los costados: un fénix cuyas alas de fuego azul parecía desplegarse hacia atrás, fundiéndose con los difusores de calor. Pero Ethan, mientras pasaba una gamuza por el asiento de polímero, sentía una inquietud que no lo dejaba descansar.

Tenía la máquina. Tenía el apoyo. Pero le faltaba el alma del piloto. Podía hacer que la moto levitara, sí; podía realizar giros básicos y aceleraciones cortas, pero en las grabaciones de las carreras profesionales veía algo diferente. Veía pilotos que no luchaban contra el aire, sino que parecían ser parte de él.

---Necesitas un maestro ---sentenció Sony una noche de junio, mientras compartían el silencio del techo del edificio de Ethan. A lo lejos, el domo de Goglee City brillaba como una joya inalcanzable, rodeada de nubes de contaminación que brillaban con un púrpura enfermizo---. Carlos sabe de tornillos, Ethan, pero no sabe qué se siente cuando la gravedad intenta partirte por la mitad.

---Papi me ha enseñado todo lo que sé ---replicó Ethan, aunque en el fondo sabía que ella tenía razón---. ¿Quién más hay en este agujero que sepa manejar una Clase 3?

Sony lo miró de reojo, su pelo rosa ondeando levemente en la brisa tóxica.

---Hay alguien. Tomás Vega. En su época lo llamaban ``El Viejo Toro''. Ganó El Velo en el 82, cuando las reglas todavía permitían un poco de honor. Es la única leyenda viva del Distrito RD, aunque ahora sea más un fantasma que un hombre.

Ethan había escuchado los susurros. Tomás había sido un dios para los mecánicos del sector. Pero después de rechazar la ciudadanía de la Ciudad Aérea como protesta por las condiciones del Distrito, las corporaciones lo borraron del mapa. Le quitaron su licencia, confiscaron su equipo y lo abandonaron a su suerte en las Ruinas Exteriores.

---Dicen que se volvió loco ---susurró Ethan---. Que habla con las máquinas muertas.

---Viene con el oficio ---dijo Sony, levantándose---. Si quieres ganar, mañana iremos a buscarlo. Pero prepárate, niño: el Toro no tiene paciencia para los soñadores que no están dispuestos a sangrar.

\bigskip
\begin{center}***\end{center}
\bigskip

Las Ruinas Exteriores eran el eco de la Gran Inundación. Allí, el mar había reclamado los niveles inferiores, dejando solo los esqueletos de acero y hormigón de la antigua ciudad. El aire era pesado, cargado de un moho milenario y el eco constante del agua golpeando los cimientos rotos.

Sony guio a Ethan hasta el sótano de lo que alguna vez fue un centro comercial de lujo. Entre escombros y cables de fibra óptica que colgaban como lianas, vivía Tomás Vega.

El hombre era más viejo de lo que Ethan imaginaba. Su piel era un pergamino de cicatrices y tatuajes borrosos de antiguas escuderías. Sus manos, nudosas y curtidas, temblaban levemente mientras ajustaba una lámpara de aceite, pero su mirada era un escáner láser que pareció analizar la estructura molecular de Ethan en un segundo.

---Sony... ---la voz de Tomás era como el crujido de piedras molidas---. Te dije que no volvieras a traerme piezas defectuosas.

---Esta no es una pieza, viejo ---respondió Sony con su habitual desparpajo cubano---. Es un piloto. O al menos, cree que lo es.

Tomás soltó una carcajada amarga que terminó en una tos seca. Se acercó a Ethan, rodeándolo como si fuera un modelo de motowave candidato al desguace.

---Diecisiete años... ---murmuró Tomás---. La edad de la estupidez infinita. ¿Qué te hace pensar que puedes sobrevivir a El Velo, chamaco? ¿Crees que porque tu moto brilla y tienes una chica guapa al lado ya eres invencible?

---No creo ser invencible ---respondió Ethan, manteniendo la espalda recta---. Pero tengo la motowave más rápida que ha visto este distrito desde que tú dejaste de correr. Y tengo el hambre que ellos, los de la ciudad alta, nunca entenderán.

---El hambre te hace descuidado ---sentenció Tomás, sentándose en un banco de madera podrida---. El Velo no es una carrera, es una trampa gravitatoria. Los de Goglee City usan los aceleradores para jugar contigo como si fueras un ratón en una licuadora. ¿Sabes lo que es ver a tu mejor amigo desintegrarse porque su sistema de fase falló por un microsegundo?

---Sé lo que es perder cosas ---dijo Ethan, pensando en su madre y en los sueños rotos de Carlos---. Por eso mismo tengo que hacerlo. No para ser un héroe, sino para que ellos dejen de ser los únicos que escriben la historia.

Tomás se quedó en silencio, mirando fijamente la llama de su lámpara. Por un momento, el brillo de sus ojos pareció viajar décadas atrás.

---Si te entreno, Ethan, no habrá piedad. Te haré odiar la gravedad antes de enseñarte a amarla. Si estás dispuesto a eso... trae tu cacharro mañana al cañón del Sector 7.

\bigskip
\begin{center}***\end{center}
\bigskip

El entrenamiento fue una pesadilla geométrica. Tomás no le enseñó a acelerar ni a frenar; lo obligó a volar a través de un laberinto de antenas caídas donde los campos magnéticos residuales hacían que la brújula de la moto girara como loca.

---¡Siente la geodésica, maldita sea! ---gritaba Tomás por la radio, su voz distorsionada por la estática---. ¡Deja de pelear con la moto! El espacio no es recto, Ethan. ¡Se curva! ¡Si luchas contra la curva, el Efecto Honda te va a escupir contra el muro!

Ethan se estrelló catorce veces el primer día. El carenado del fénix, tan cuidadosamente pintado por Yeisly, ahora tenía cicatrices de concreto y metal. Al final de la tarde, Ethan apenas podía sostener el manillar; sus músculos gritaban de fatiga y su mente era un caos de vectores y fuerzas G.

---No puedo... no veo lo que quieres que vea ---jadeó Ethan, desplomándose sobre el suelo de ceniza.

Tomás se acercó lentamente, cojeando de su pierna cibernética mal ajustada. Se sentó a su lado y dibujó un arco en el polvo.

---Imagínate que el espacio-tiempo es una sábana tensa, Ethan. Si pones un objeto pesado, la sábana se hunde. El camino más corto entre dos puntos en ese hundimiento no es una línea recta. Es una línea curva... una geodésica. La mayoría de los pilotos usan sus motores para forzar la recta, gastando energía y perdiendo estabilidad. Pero los grandes... los que de verdad vuelan... ellos ven el hundimiento. Ellos surfean la curva.

---Luna... ---susurró Ethan, recordando el nombre que Sony había mencionado.

El rostro de Tomás se ensombreció. Sacó un colgante de su cuello, una pequeña placa de identificación de piloto, rayada y vieja.

---Luna Reyes era mi compañera. Ella veía las geodésicas mejor que nadie. Podía entrar en un flujo gravitatorio y desaparecer, ganando segundos imposibles. Pero el Surfeo de Geodésicas tiene un precio: si tu ángulo de entrada falla por tres milímetros, la curvatura se vuelve en tu contra. El espacio te expulsa como una honda a ochocientos kilómetros por hora. No hay frenos, no hay control. Solo el impacto.

---Por eso murió... --- Ethan miró la placa.

---Murió buscando la geodésica perfecta. ---Tomás guardó la placa---. El Velo pone Martillos Gravitacionales en los puntos de salida de las curvas. Son muros de gravedad pura. Si te expulsa el Efecto Honda, te conviertes en una pintura de sangre en la pantalla de los patrocinadores.

Ethan cerró los ojos y trató de proyectar mentalmente el cañón. Visualizó los campos de fuerza, los hundimientos, las tensiones invisibles del aire. De repente, el zumbido de su propia moto empezó a tener sentido. No era un ruido mecánico; era una respuesta a la frecuencia del mundo.

---Lo siento ---susurró Ethan---. El pulso. Puedo sentir el pulso del motor cambiando cuando paso cerca del edificio.

Tomás se detuvo en seco. Sus ojos se abrieron con una chispa de esperanza que intentó ocultar rápidamente.

---¿Sientes el cambio de frecuencia?

---Sí. Es como si la moto supiera por dónde quiere ir.

Tomás apretó los puños. Sabía que eso no se podía enseñar. Era el Don. El mismo sentido extrasensorial que había hecho de Luna una leyenda y, finalmente, una mártir.

---Si puedes sentir eso, Ethan... entonces tal vez haya una oportunidad. Pero vas a necesitar algo más que instinto. Vas a necesitar la Biblia del Universo de Tomás Vega. ---Sacó una tableta antigua, protegida por una carcasa de titanio---. Aquí están mapeadas todas las geodésicas estables de los circuitos históricos. Si logras sincronizar tu instinto con esta geometría... serás imparable.

Ethan tomó la tableta con reverencia. Era el legado prohibido del hombre que había ganado El Velo y se había negado a vender su alma.

---Mañana empezaremos de verdad ---dijo Tomás, y esta vez no hubo burla en su voz---. Vamos a enseñarle a esa ciudad alta lo que pasa cuando un chico del Distrito decide que ya no quiere caminar por el suelo.

\bigskip
\begin{center}***\end{center}
\bigskip

\textbf{FIN DEL CAPÍTULO 3}


\bigskip
\begin{center}***\end{center}
\bigskip


\part{EL VELO}

\chapter{EL PESO DE LA GRAVEDAD}

\section*{Diciembre de 2114}

Otro día en el Distrito RD. La gravedad aquí no se siente solo como una fuerza física, sino como una condena social.

El cielo no es azul; hace décadas que el azul es un color de catálogo para los que viven dentro del domo. Aquí fuera, es una losa sólida e interminable de smog espeso que apenas deja pasar la luz del sol. El astro rey es un disco rojizo y enfermo, una pupila inflamada que observa un mundo de óxido y desesperanza desde detrás de una cortina de azufre. El aire no se respira, se mastica. Tiene un sabor amargo a ciclos de vida agotados, a combustible fundido y al aliento metálico de las megas-fábricas que nunca se detienen.

Las calles del Distrito RD son un laberinto de neón parpadeante que lucha contra la calima industrial. Los charcos de agua ácida en el suelo reflejan los carteles de propaganda de Goglee City, donde modelos impecables prometen ``Pureza y Paz'' bajo cielos que Ethan solo ha visto en sueños. La gente camina rápido, con la cabeza gacha, protegida por máscaras de filtro barato que suelen fallar cuando el viento sopla desde el Sector Químico.

Son los hijos de la diáspora, los nietos de la Gran Inundación. El Distrito RD no se llamaba así por casualidad; se había convertido en el último refugio para los que huyeron cuando el Atlántico se tragó las costas del Caribe. Ethan recordaba las historias de su abuelo sobre una Miami que ahora era un arrecife de coral artificial bajo cien metros de agua, de una Venezuela donde las ciudades costeras se convirtieron en archipiélagos de miseria, y de una Colombia cuyas playas fueron devoradas por un mar que no pedía perdón. Esa mezcla de sangres ---dominicana, cubana, venezolana, colombiana--- latía bajo el ruido de las prensas hidráulicas, una cultura de sol y música que ahora se veía forzada a sobrevivir en un invierno industrial permanente.

En su taller, Ethan no permitía que la atmósfera opresiva entrara en su zona de trabajo. \textit{El Fénix} descansaba en el centro del taller, ya no como la promesa oxidada de su cumpleaños, sino como una depredadora de metal pulido. Bajo la luz blanca de las lámparas LED, el chasis brillaba con una intensidad metálica que parecía absorber el humo del taller.

---Ajuste del bus de datos al 99.8 por ciento ---susurró Ethan, su voz suave pero cargada de una tensión eléctrica. Sus dedos, callosos pero precisos, ajustaban un micro-conector en el núcleo gravitacional---. Necesitas aguantar, pequeña. Los circuitos de El Velo no perdonan. Si el sistema de fase flaquea bajo la presión de un Martillo Gravitacional, no quedarán ni los tornillos para el entierro.

Rafa estaba casi sumergido dentro de la estructura de la moto, con los pies colgando sobre el motor de inducción. Su linterna frontal bailaba sobre una red de fibra óptica que emitía un suave pulso verde.

---Tranquilo manín, que este cableado lo hice yo pensando en la guerra ---dijo Rafa, emergiendo con una mancha de grasa que le cruzaba el pómulo como pintura de combate---. He reforzado los blindajes electromagnéticos. Si una onda de interferencia corporativa nos pega, la moto ni se va a enterar. Pero Ethan... el problema no es mi cobre, el problema es que el estabilizador que tenemos no va a aguantar una aceleración geodésica.

Ethan pasó la mano por el carenado frío. Sabía que Rafa tenía razón. Tomás lo había advertido: sin un estabilizador T-9, intentar un Surfeo de Geodésicas era como lanzarse a un huracán con un paraguas roto.

---El plazo de inscripción cierra en dos días ---dijo una voz desde las sombras. Carlos entró, tosiendo levemente, con un trapo impregnado de solvente en la mano. Se detuvo a observar la máquina---. ¿De verdad vas a ir a la clasificación mañana sin el T-9?

---No tenemos otra opción, papi ---respondió Ethan, aunque se le notaba la frustración en el tono---. Un T-9 cuesta más créditos de los que hemos visto juntos en toda la vida.

---La fe no arregla la física, mijo ---sentenció Carlos, acercándose y revisando el dial de presión gravitacional---. He oído que El Tocayo recibió una caja de ``excedentes militares'' del último transporte que falló en el domo. Dicen que hay tecnología de grado táctico. Si quieres una oportunidad real, deja de pulir el metal y ve a la Plaza. Ahora.

\bigskip
\begin{center}***\end{center}
\bigskip

La Plaza de Chatarra bajo la lluvia ácida era un infierno de vapor y estridencia. Ethan se movía entre los puestos, esquivando a los chatarreros que intentaban venderle filtros de aire usados y baterías descargadas. Su mente estaba fija en el puesto de El Tocayo.

Cuando llegó, el viejo estaba sentado sobre una caja de municiones vacía, con su mandíbula de cromo chirriando mientras comía algo que Ethan prefirió no identificar.

---Vienes por el T-9, ¿verdad? ---gruñó El Tocayo antes de que Ethan dijera una palabra---. La noticia corre rápido en la basura.

---¿Lo tienes? ---Ethan trató de que su voz no temblara.

El Tocayo escupió a un lado y pateó una lona negra. Debajo, descansaba una pieza de geometría compleja, con aletas de refrigeración de cerámica y un puerto de conexión dorado. Un estabilizador de fase T-9. Grado militar.

---Es una belleza de la vieja Guardia ---dijo El Tocayo---. Quinientos créditos. Y no acepto promesas, solo créditos fríos.

---Casi no tenemos esa cantidad ---dijo una voz a la espalda de Ethan.

Sony estaba allí, protegida por un chubasquero de plástico que goteaba sobre el pavimento sucio. Se acercó a la pieza y la analizó con una mirada cargada de un cinismo experto.

---Oye, El Tocayo, no te pases de listo con el chico ---soltó Sony, su acento cubano cortando el ruido de la lluvia---. Esa pieza tiene los sellos de seguridad violados. Es material robado o ``caído del camión''. Si los recolectores corporativos la ven aquí, te queman el puesto.

El Tocayo gruñó, pero no se inmutó.

---La seguridad cuesta, niña revolucionaria. Si lo quiere, que pague.

---Dáselo por cuatrocientos ---negoció Sony, poniendo una mano sobre el hombro de Ethan---. Él es el único que puede hacer que esa pieza valga algo más que su peso en metal. Si él gana, todos los ojos estarán en el Distrito. ¿Acaso no quieres ver a los popis llorar cuando un fénix de chatarra los adelante por la izquierda?

\bigskip
\begin{center}***\end{center}
\bigskip

De vuelta en el santuario del taller, el silencio era casi absoluto, solo roto por el zumbido de la estática de la radio de Rafa. Ethan trabajaba con una intensidad febril, instalando el T-9 bajo la supervisión técnica de Unica, quien le dictaba los voltajes de calibración desde un monitor auxiliar.

---Voltaje de pre-mantenimiento estable, Ethan ---dijo Unica, su voz educada pero cargada de una extraña emoción contenida---. La resonancia de fase está en verde. Puedes iniciar el núcleo.

Ethan activó la secuencia. Esta vez, el sonido de \textit{El Fénix} era diferente. Ya no era un zumbido, era un ronroneo profundo, una vibración subsónica que hacía que los pelos de la nuca se erizaran. La moto se elevó, pero esta vez lo hizo con una gracia absoluta, sin la más mínima oscilación.

---¡Mierda, manito, parece que flota en seda! ---exclamó Rafa, asombrado por la estabilidad.

Ethan se acercó a su padre, quien observaba la máquina desde el fondo del taller. Se arrodilló y tomó la mano callosa de Carlos, llevándosela a la frente en el ritual que nunca olvidaba.

---Ción, papi. Mañana empieza todo.

Carlos puso su mano pesada y cálida sobre la cabeza de su hijo. Su voz era un susurro que luchaba contra la tos persistente.

---Dios te me bendiga, mijo. No olvides lo que te enseñó Tomás. La gravedad es el camino, no el enemigo. Vuela tan alto que ellos no puedan ni ver tu rastro de ozono.

A través del tragaluz sucio del taller, Ethan miró hacia arriba. El domo de Goglee City brillaba como una estrella inalcanzable de cristal y privilegios. Mañana, por primera vez, él no sería solo un espectador desde el fango. Mañana, él sería la fuerza que desafiaría el orden del mundo.

El motor de \textit{El Fénix} rugió, un grito de desafío que resonó en las paredes de chapa y óxido. Ethan Ramírez estaba listo para el Velo.

\bigskip
\begin{center}***\end{center}
\bigskip

\textbf{FIN DEL CAPÍTULO 4}


\bigskip
\begin{center}***\end{center}
\bigskip

\chapter{CHISPAS EN EL ÓXIDO}

Los dos días siguientes fueron un borrón de grasa, sudor y el chirrido constante de las herramientas neumáticas. Ethan apenas durmió; cada vez que cerraba los ojos, veía líneas de geodésicas trazadas en el interior de sus párpados. \textit{El Fénix} necesitaba ser más que una moto; necesitaba ser una parte de su sistema nervioso. En las afueras del Distrito RD, donde las ruinas de los antiguos rascacielos formaban cañones de metal oxidado que desafiaban la física, Ethan se dedicó a romper sus propios límites.

El sol enfermo se asomaba entre el smog, tiñendo el mundo de un naranja sucio y radiactivo. Ethan aceleró por una avenida que el mar había dejado seca hace una década, sintiendo el zumbido de los propulsores bajo sus piernas como el ronroneo de una bestia impaciente.

---¡Estabilizadores de fase al 100%! ---gritó por el intercomunicador, su voz compitiendo con el aullido del viento sintético---. ¡Mantenme íntegro, pequeña!

El T-9, la pieza militar que Sony había conseguido en la Plaza, vibraba con un calor reconfortante bajo el asiento. Ethan sentía la frecuencia del estabilizador sincronizándose con su propio pulso; era una coraza invisible que evitaba que la inercia lo hiciera pedazos en cada giro. Pero la estabilidad era solo la base; para ganar necesitaba la magia de Unika.

---¡Unika, necesito el desfase ahora! ---exclamó Ethan, inclinando la moto hacia una rampa de escombros.

---No me grites, Ethan, los sensores están funcionando perfectamente ---la voz de Unika llegó nítida, precisa y con esa calma técnica que siempre lograba bajarle las revoluciones al piloto---. Inversor G activo. La polaridad local está invertida. Tienes cinco segundos de tracción vertical antes de que el núcleo necesite recargarse. Úsalos bien.

---¡Te copio, manín! ---la voz de Rafa entró en la frecuencia, rebosante de entusiasmo---. ¡Dale con todo! He puenteado los relés de salida para que no te den tirones. ¡Si esa moto fuera una mujer, tendría que casarme con ella de lo bien que suena!

La motowave rugió y, en un parpadeo, las ruedas se plegaron con un \textit{clic} hidráulico. La levitación se activó con una agresividad nueva. Ethan salió disparado hacia arriba, corriendo literalmente por la pared de cristal roto de un edificio de cincuenta pisos. La gravedad intentaba arrancarlo de la superficie, pero el Inversor G de Unika mantenía la adherencia molecular. Para Ethan, en ese instante, el mundo se había girado noventa grados.

Al llegar a la azotea, Ethan se lanzó al vacío, realizando un giro de trescientos sesenta grados que lo dejó de frente a la calle inferior. Aterriza con una suavidad imposible, amortiguado por un colchón de aire gravitatorio.

Se quitó el casco, jadeando, mientras el sudor le corría por la frente. Unika se acercó desde la sombra de un hangar abandonado. Caminaba con su habitual elegancia austera, sosteniendo la tableta de diagnósticos como si fuera un escudo.

---Los datos de telemetría indican que todavía eres demasiado brusco en la entrada de la geodésica ---dijo ella, sin levantar la vista de la pantalla, aunque Ethan pudo notar un leve temblor en su voz---. Si el Inversor G detecta una fluctuación de más de 0.5 milisegundos, la moto recuperará su peso real y caerás como una piedra.

---Pero no caí ---sonrió Ethan, bajándose de la moto---. Porque tú construiste el mejor inversor de este lado del domo. ¿Dónde aprendiste a programar así, Unika?

Ella se detuvo y, por primera vez, hubo un silencio largo. Miró hacia el domo brillante de Goglee City, con una expresión que Ethan nunca había visto: una mezcla de nostalgia y rencor.

---Mis padres eran ingenieros de sistemas en la Ciudad Aérea ---confesó ella en un susurro---. Yo nací allí, Ethan. Vivíamos en el Nivel 4. Pero cometieron el ``error'' de detectar un desvío de fondos corporativos hacia los proyectos militares del Velo. Una noche, simplemente perdimos la ciudadanía. Nos ``deseleccionaron''. Nos enviaron aquí abajo con nada más que la ropa que llevábamos puesta.

Ethan se quedó helado. Nunca habría imaginado que Unika conocía el interior de las murallas.

---Por eso eres tan perfeccionista ---murmuró él---. Quieres volver.

Unika lo miró fijamente con sus ojos oscuros, llenos de una determinación fría.

---No quiero volver para vivir como ellos. Quiero que sientan lo que es perderlo todo ante alguien que ellos consideran ``desecho''. Por eso el Inversor G es perfecto. Porque es mi venganza empaquetada en un chip.

De repente, Yeisly apareció por la esquina del taller, interrumpiendo el momento con una risotada.

---¡Muchachos! ¡Dejen la charla que ya tengo el diseño final! ---Yeisly traía una lata de pintura especial de secado rápido---. Primo, le puse unos detalles en neón reactivo a los costados. Cuando aceleres, el fénix va a parecer que está quemando el aire. ¡Qué maldita vaina más jebi se va a ver!

Rafa llegó detrás, cargando una caja de herramientas.

---Manito, ya estamos. El Fénix está listo para el Velo. Mañana, el Distrito RD va a tener un piloto en la pista oficial.

\bigskip
\begin{center}***\end{center}
\bigskip

La mañana de la carrera llegó con una tormenta eléctrica que parecía un presagio. El cielo retumbaba con una furia metálica y los relámpagos iluminaban el smog como venas de luz purpúrea. Ethan se despidió de Carlos en la entrada del taller. 

Su padre estaba más pálido de lo normal, pero se mantenía firme, con la espalda recta a pesar de la respiración dificultosa que revelaba su enfermedad pulmonar avanzada.

---Escúchame, mijo ---dijo Carlos, poniendo sus manos pesadas sobre los hombros de Ethan---. Los pilotos de Goglee City no son deportistas; son asesinos con trajes caros. Tienen órdenes de eliminar a cualquier independiente que amenace su ``espectáculo''. Los Martillos Gravitacionales que verás en la pista son trampas diseñadas para los que no tienen equipo de fábrica.

---Tengo a Unika, a Rafa y a ti ---respondió Ethan con firmeza---. Eso es más de lo que ellos tendrán nunca.

---Confía en lo que Tomás te grabó en la mente ---insistió Carlos---. La gravedad no es algo que se vence, es algo con lo que se baila.

Ethan se arrodilló, tomó la mano derecha de su padre y la llevó a su frente.

---Ción, papi.

---Dios te bendiga, hijo mío. No me traigas un trofeo. Tráeme el orgullo de saber que sigues vivo.

\bigskip
\begin{center}***\end{center}
\bigskip

El viaje hacia la frontera del domo fue una bofetada sensorial. A medida que se acercaban, el aire se volvía más puro, libre de las partículas de carbón y azufre del Distrito. Para Ethan, respirar ese aire era casi doloroso; sus pulmones gritaban ante la falta de impurezas a las que se habían adaptado.

La Puerta del Velo era un arco colosal de luz blanca. Un escáner biométrico revisó a Ethan mientras drones de seguridad orbitaban a su alrededor como avispas mecánicas. Al otro lado, la opulencia: boxes de cristal templado, mecánicos con uniformes impecables y motowaves que parecían diseñadas por joyeros.

Un corredor con el logo de ``Hyperion Dynamics'' se detuvo junto a Ethan. Su casco reflejaba el entorno con un acabado de espejo.

---¿De qué vertedero sacaste ese juguete, rata del Distrito? ---se burló el corredor a través del canal de proximidad abierto.

Ethan no respondió. Simplemente miró a Unika, que estaba al borde de la pista con su tableta lista. Ella asintió. Ethan activó el Inversor G y el chasis de \textit{El Fénix} vibró con un zumbido agudo y amenazante que silenció las risas de los mecánicos cercanos.

---¡Pilotos a sus posiciones! ---la voz de la inteligencia artificial de la carrera resonó en todo el estadio---. ¡Preparando secuencia de inversión gravitacional!

Ethan bajó el visor de su casco. El mundo se volvió un mapa de vectores de luz azul. Buscó la geodésica inicial, el camino invisible que solo él podía ver. El corazón le latía con fuerza contra las costillas, cada latido sincronizado con el núcleo de la moto.

---\textit{En 3... 2... 1... ¡IGNICIÓN!}

El mundo se invirtió. Ethan se lanzó al vacío, no como una víctima de la gravedad, sino como su dueño.

\bigskip
\begin{center}***\end{center}
\bigskip

\textbf{FIN DEL CAPÍTULO 5}


\bigskip
\begin{center}***\end{center}
\bigskip

\chapter{EL VELO GRAVITACIONAL}

La velocidad no es una cifra en un panel holográfico. Es una distorsión de la realidad.

A seiscientos kilómetros por hora, el Distrito RD se convirtió en un borrón de tonos sepias y grises que se desvanecía en el espejo retrovisor. Ethan y \textit{El Fénix} eran un relámpago azul cobalto surcando una pista de cristal sintético que parecía flotar sobre el abismo de las ruinas. A su alrededor, la competencia no eran simples pilotos; eran depredadores corporativos en máquinas que valían más que todo un sector del Distrito.

La pista del Velo era una serpiente caprichosa que se retorcía en el aire, desafiando toda lógica euclidiana. La gravedad cambiaba cada tres segundos, orquestada por los generadores de Goglee City. Ethan sentía cómo sus órganos se desplazaban con cada giro; arriba se convertía en un suelo sólido, y la izquierda era un muro que lo succionaba hacia el vacío.

---¡Calibración de fase al límite, Ethan! ---la voz de Unika sonó en su casco, tensa pero precisa---. El núcleo está operando al 115 por ciento. No le pidas más o el inversor estallará.

---¡Dale, manito, que ese Viper te tiene en la mira! ---Rafa gritaba desde el taller clandestino, donde Yeisly y él monitoreaban la frecuencia pirateada---. ¡Cuidado con la maniobra de cierre, que ese tipo juega puerco!

Viper, el piloto estrella de Hyperion Dynamics, era una sombra negra y dorada que se movía con una fluidez insultante. Embistió el flanco derecho de \textit{El Fénix}, un impacto de metal contra metal que hizo saltar chispas de titanio y vibrar hasta el último tornillo de la moto artesanal. Pero la máquina de Ethan no estaba hecha de polímeros delicados de laboratorio; estaba forjada con la terquedad de los que no tienen nada que perder.

---¡Martillos Gravitacionales activos en el Sector 4! ---advirtió la IA de la pista.

Del cielo descendieron columnas de fuerza pura, pistones invisibles que golpeaban el asfalto con la potencia de meteoritos. Ethan vio con horror cómo un corredor de Neos-Tech fallaba en la sincronización; su vehículo fue aplastado en medio segundo, convirtiéndose en una tortilla de metal humeante que la pista escupió hacia las ruinas inferiores.

\textit{``No luchas contra el océano, hijo. Te conviertes en la ola.''}

La voz de Carlos resonó bajo el ruido ensordecedor de los propulsores. Ethan inhaló el poco oxígeno que le quedaba y tomó una decisión suicida.

---Unika, voy a intentar el Surfeo de Geodésicas en la curva del Martillo.

---¡Es un suicidio! ---gritó Unika---. El ángulo de entrada es demasiado estrecho. Si fallas por tres milímetros, el Efecto Honda te lanzará al espacio exterior.

---Es el único camino. ¡Rafa, abre el flujo de energía trasera!

Ethan giró el manillar violentamente y salió de la pista principal, lanzándose hacia la barandilla exterior. Por un instante, \textit{El Fénix} voló sobre el vacío total. El Efecto Honda lo golpeó; sintió cómo la gravedad invertida intentaba lanzarlo hacia atrás como un proyectil fuera de control. El chasis chirrió, los cables de Rafa empezaron a humear y el Inversor G de Unika emitió un grito sordo de sobrecarga.

Entonces, el tiempo pareció detenerse. Ethan cerró los ojos, no por miedo, sino para \textit{ver}. Recordó la técnica de Tomás, el ``Viejo Toro''. Las líneas invisibles de la geodésica aparecieron en su mente como hilos de luz dorada.

---Te tengo ---susurró Ethan.

Encontró el valle gravitatorio. \textit{El Fénix} se enganchó a la curvatura del espacio-tiempo. Donde antes había resistencia, ahora había un impulso infinito. La moto se disparó hacia adelante con una aceleración que rompió la barrera del sonido en una explosión de luz azul y vapores de ozono.

Viper fue un borrón que Ethan dejó atrás en una décima de segundo. El piloto corporativo intentó reaccionar, pero la física ya no estaba de su lado.

Ethan cruzó la línea de meta entrando en el Domo.

El cambio fue instantáneo. Del ruido infernal de los motores y el aire turbulento, pasó a un silencio absoluto y una paz artificial. \textit{El Fénix} se detuvo con elegancia frente a la tribuna principal de Goglee City, donde los hombres y mujeres más poderosos del mundo se quedaron en un silencio sepulcral.

Ethan Ramírez, el chico que reparaba chatarra en un taller con techos de chapa, acababa de ganar la carrera más prestigiosa del planeta.

\bigskip
\begin{center}***\end{center}
\bigskip

Ethan se quitó el casco. Sus manos temblaban, no de fatiga, sino de una adrenalina que todavía no encontraba salida. Respiró el aire de Goglee City. Era un aire sin alma, sin olor a mar ni a combustible; un aire esterilizado que sabía a hospital.

Un funcionario corporativo de alto rango, con un traje de seda gris que parecía cambiar de color con la luz, se acercó flotando sobre una plataforma de levitación personal. Traía consigo un pase holográfico que emitía un brillo dorado, el símbolo de la Ciudadanía Aérea.

---Felicidades, Ciudadano Ramírez ---dijo el funcionario con una sonrisa ensayada que no llegaba a sus ojos---. Has demostrado una aptitud excepcional. Bienvenido a la verdadera civilización. Tu vida en el barro ha terminado hoy.

Ethan miró el pase. Era el sueño de Carlos. Era la seguridad para Unika. Era el fin del hambre para Rafa y Yeisly. Pero luego miró a través de los muros de cristal reforzado del domo. Allá abajo, el Distrito RD seguía oculto bajo la calima tóxica. Allá abajo, el fuego de la gente seguía ardiendo en medio de la oscuridad.

---¿Y qué pasa con mi gente? ---preguntó Ethan, su voz ronca por el esfuerzo pero clara como el cristal.

---El premio es individual, muchacho ---respondió el funcionario, restándole importancia con un gesto---. Los desechos no entran en el centro de procesamiento. Ahora eres uno de nosotros.

Ethan apretó el pase en su puño y luego, ante la mirada atónita de las cámaras que transmitían a todo el mundo, lo dejó caer al suelo sintético.

---Guárdese su paraíso de plástico ---sentenció Ethan, poniéndose de pie sobre el asiento de su moto---. No gané para ser uno de ustedes. Gané para demostrar que ustedes no son mejores que nosotros.

Un murmullo de horror recorrió la tribuna. En el Distrito RD, el grito de júbilo fue tan fuerte que hizo vibrar los cimientos de las ruinas.

---Si no hay lugar para mi gente aquí arriba ---continuó Ethan, mirando fijamente a la cámara---, entonces bajaremos el domo al suelo.

Ethan se puso el casco. El motor de \textit{El Fénix} rugió una vez más, no como un transporte, sino como una declaración de guerra. Giró la moto y salió de la zona de boxes, no hacia su nuevo apartamento de lujo, sino de vuelta hacia el túnel de salida.

Volvió afuera. Pero esta vez, no volvió como una rata. Volvió como la chispa que incendiaría el velo.

\bigskip
\begin{center}***\end{center}
\bigskip

En la Plaza de Chatarra, Sony contemplaba la pantalla con lágrimas de orgullo y una sonrisa que prometía cambios drásticos.

---Eso es, Ethan... ---susurró Sony, ajustando su mochila llena de explosivos electrónicos---. Esa es la chispa. Ahora, vamos a prender el fuego.

En el taller, Carlos cerró los ojos y respiró hondo, sintiendo por primera vez que el aire tóxico olía un poco más a esperanza.

---Dios te bendiga, hijo mío. Has roto la gravedad.

Ethan Ramírez no solo había ganado una carrera. Había empezado una revolución.

\bigskip
\begin{center}***\end{center}
\bigskip

\textbf{FIN DEL CAPÍTULO 6}

\end{document}
